\documentclass[11pt]{beamer}

% Hỗ trợ tiếng Việt
\usepackage[utf8]{vietnam}
\usepackage{graphicx}
\usepackage{tikz} % Để vẽ đồ thị step-by-step
\usetikzlibrary{positioning, arrows.meta}

% Giao diện slide (bạn có thể đổi Madrid thành các theme khác như: default, Boadilla, CambridgeUS...)
\usetheme{Madrid}
\usecolortheme{default}

% ==========================================
% ĐOẠN CODE THÊM VÀO ĐỂ ĐỔI MÀU & CHÂN TRANG (GIỮ NGUYÊN NỘI DUNG BÊN DƯỚI)
% ==========================================
\definecolor{DarkTeal}{RGB}{0, 80, 80}
\definecolor{OrangeAccent}{RGB}{255, 150, 0}

\setbeamercolor{structure}{fg=DarkTeal}
\setbeamercolor{palette primary}{bg=DarkTeal, fg=white}
\setbeamercolor{palette secondary}{bg=OrangeAccent, fg=white}
\setbeamercolor{palette tertiary}{bg=DarkTeal!90, fg=white}
\setbeamercolor{frametitle}{bg=DarkTeal, fg=white}
\setbeamercolor{block title}{bg=DarkTeal, fg=white}
\setbeamertemplate{navigation symbols}{} % Ẩn thanh điều hướng

% Tùy chỉnh chân trang giống ảnh
\setbeamertemplate{footline}{
  \leavevmode%
  \hbox{%
  \begin{beamercolorbox}[wd=.333333\paperwidth,ht=2.25ex,dp=1ex,center]{palette tertiary}%
    \usebeamerfont{author in head/foot}\insertshortauthor
  \end{beamercolorbox}%
  \begin{beamercolorbox}[wd=.333333\paperwidth,ht=2.25ex,dp=1ex,center]{palette secondary}%
    \usebeamerfont{title in head/foot}LÝ THUYẾT ĐỒ THỊ
  \end{beamercolorbox}%
  \begin{beamercolorbox}[wd=.333333\paperwidth,ht=2.25ex,dp=1ex,right]{palette primary}%
    Học kỳ II - 2026 \hspace*{2em}
    \insertframenumber{} / \inserttotalframenumber \hspace*{2ex} 
  \end{beamercolorbox}}%
  \vskip0pt%
}
% ==========================================

% Cài đặt thông tin trang bìa
\title[Lý Thuyết Đồ Thị]{Bài Thuyết Trình: Lý Thuyết Đồ Thị \& Ứng Dụng}
\subtitle{Phần 3: Thuật toán \& Phần 4: Các bài toán hay}
\author[Tung Tung Tung Sahur]{Tung Tung Tung Sahur} % Đã cập nhật tên bạn
\date{\today}

\begin{document}

% --- TRANG BÌA ---
\begin{frame}
    \titlepage
\end{frame}

% --- MỤC LỤC ---
\begin{frame}{Nội dung chính}
    \tableofcontents
\end{frame}

% ==========================================
% PHẦN 3: THUẬT TOÁN
% ==========================================
\section{Phần 3: Thuật toán}

\begin{frame}{ PHẦN 3: Thuật toán}
    \begin{center}
    \end{center}
    \vspace{0.5cm}
    \begin{itemize}
        \item[$\bullet$] \textbf{BFS} (Breadth-First Search - Tìm kiếm theo chiều rộng)
        \begin{itemize}
        \item \textit{Trình bày có hình minh họa từng bước (step-by-step).}
        \end{itemize}
        \vspace{0.3cm}
        \item[$\bullet$] \textbf{DFS} (Depth-First Search - Tìm kiếm theo chiều sâu)
        \vspace{0.3cm}
        \item[$\bullet$] \textbf{Dijkstra}
        \begin{itemize}
        \item \textit{Giới thiệu cơ bản về thuật toán tìm đường đi ngắn nhất.}
        \end{itemize}
    \end{itemize}
    
    \vspace{0.5cm}
    \begin{block}{Ghi chú}
         Tất cả các thuật toán đều được làm theo dạng \textbf{"step-by-step"}.
    \end{block}
\end{frame}

% --- BFS STEP BY STEP (Ví dụ minh họa) ---
\begin{frame}{1. Thuật toán BFS (Tìm kiếm theo chiều rộng)}
    \textbf{Minh họa Step-by-step:} Bắt đầu từ đỉnh A.
    \vspace{0.5cm}
    
    \begin{center}
    \begin{tikzpicture}[
        node distance=2cm,
        main node/.style={circle, draw, font=\bfseries, minimum size=8mm},
        visited/.style={circle, draw, fill=blue!30, font=\bfseries, minimum size=8mm},
        queue/.style={circle, draw, fill=green!30, font=\bfseries, minimum size=8mm}
    ]
        % Khai báo các node
        \node[main node]<1>[label=above:Start] (A) {A};
        \node[visited]<2->[label=above:Start] (A) {A};
        
        \node[main node]<1-2> (B) [below left of=A] {B};
        \node[queue]<3> (B) [below left of=A] {B};
        \node[visited]<4-> (B) [below left of=A] {B};
        
        \node[main node]<1-2> (C) [below right of=A] {C};
        \node[queue]<3-4> (C) [below right of=A] {C};
\node[visited]<5-> (C) [below right of=A] {C};
        
        \node[main node]<1-4> (D) [below of=B] {D};
        \node[queue]<5> (D) [below of=B] {D};
        \node[visited]<6-> (D) [below of=B] {D};

        % Khai báo các cạnh
        \path[draw, thick]
        (A) edge (B)
        (A) edge (C)
        (B) edge (D)
        (C) edge (D);
    \end{tikzpicture}
    \end{center}
    
    \vspace{0.5cm}
    \begin{itemize}
        \item<1-> \textbf{Bước 0:} Đồ thị ban đầu.
        \item<2-> \textbf{Bước 1:} Thăm đỉnh A, đưa vào hàng đợi.
        \item<3-> \textbf{Bước 2:} Lấy A ra, đưa các đỉnh kề của A (B, C) vào hàng đợi.
        \item<4-> \textbf{Bước 3:} Lấy B ra, đưa đỉnh kề (D) vào hàng đợi.
        \item<5-> \textbf{Bước 4:} Lấy C ra (không có đỉnh kề mới).
        \item<6-> \textbf{Bước 5:} Thăm D. Hoàn thành BFS.
    \end{itemize}
\end{frame}

% --- DFS ---
\begin{frame}{2. Thuật toán DFS (Tìm kiếm theo chiều sâu)}
    \begin{itemize}
        \item \textbf{Khái niệm:} Khám phá các nhánh sâu nhất có thể trước khi quay lui (backtracking). \pause
        \item \textbf{Cấu trúc dữ liệu:} Sử dụng \textit{Stack (Ngăn xếp)}. \pause
        \item \textbf{Quy trình cơ bản:}
        \begin{enumerate}
        \item Thăm một đỉnh kề chưa được thăm. Đánh dấu đã thăm và đẩy vào Stack. \pause
        \item Nếu không tìm thấy đỉnh kề nào, lấy đỉnh ở trên cùng của Stack ra. \pause
        \item Lặp lại đến khi Stack rỗng.
        \end{enumerate}
    \end{itemize}
    \vspace{0.5cm}
    \textit{(Bạn có thể copy đoạn code vẽ sơ đồ TikZ ở slide BFS trên, đổi màu để minh họa cách đi sâu vào 1 nhánh của DFS cho các slide tiếp theo).}
\end{frame}

% --- DIJKSTRA ---
\begin{frame}{3. Thuật toán Dijkstra (Cơ bản)}
    \begin{block}{Mục đích}
        Tìm đường đi ngắn nhất từ một đỉnh nguồn đến tất cả các đỉnh còn lại trong đồ thị (có trọng số không âm).
    \end{block}
    \pause
    
    \textbf{Ý tưởng Step-by-step:}
    \begin{itemize}
        \item<2-> Khởi tạo khoảng cách từ điểm đầu đến chính nó là $0$, đến các điểm khác là $\infty$.
        \item<3-> Chọn đỉnh có khoảng cách nhỏ nhất chưa được duyệt.
        \item<4-> Cập nhật lại khoảng cách cho các đỉnh kề của nó.
        \item<5-> Lặp lại cho đến khi tất cả các đỉnh đều được duyệt.
    \end{itemize}
\end{frame}

\end{document}
